\documentclass[a4paper]{article}
\usepackage{times}
\usepackage[utf8]{inputenc}
\usepackage[T1]{fontenc}
\usepackage{selinput}
\usepackage{upquote}
\usepackage[margin=2cm, rmargin=4cm, tmargin=3cm]{geometry}
\usepackage{tcolorbox}
\tcbuselibrary{breakable}
\usepackage{xspace}
\usepackage{amsmath}
\usepackage{float}
\usepackage[french]{babel}
\usepackage{hyperref}
\usepackage{xurl}
\usepackage{fontawesome6}
\usepackage{enumitem}
\usepackage{marginnote}
\usepackage{ulem}
\usepackage[yyyymmdd,hhmmss]{datetime}
\usepackage{ulem}
\usepackage{environ}
\usepackage{graphicx}
%\usepackage[top=Bcm, bottom=Hcm, outer=Ccm, inner=Acm, heightrounded, marginparwidth=Ecm, marginparsep=Dcm]{geometry}
\usepackage{courier}



\newtcolorbox{Example}[1]{colback=white,left=20pt,colframe=slideblue,fonttitle=\bfseries,title=#1}
\newtcolorbox{Solutions}[1]{breakable,colback=white,left=20pt,colframe=green,fonttitle=\bfseries,title=#1}
\newtcolorbox{Conseils}[1]{colback=white,left=20pt,colframe=slideblue,fonttitle=\bfseries,title=#1}
\newtcolorbox{Warning}[1]{colback=white,left=20pt,colframe=warning,fonttitle=\bfseries,title=#1}
\newtcolorbox{Note}[1]{colback=white,left=20pt,colframe=grey,fonttitle=\bfseries,title=#1}
\newtcolorbox{Dessin}[1]{colback=white,left=20pt,colframe=grey,fonttitle=\bfseries,title=#1}

\setlength\parindent{0pt}

  %Exercice environment
  \newcounter{exercice}
  \newenvironment{Exercice}[1][]
  {
  \par
  \stepcounter{exercice}\textbf{Question \arabic{exercice}:} (\faClock \enskip \textit{#1})
  }
  {\bigskip}
  

% Title
\newcommand{\titre}{\begin{center}
  \section*{Algorithmes et Pensée Computationnelle}
\end{center}}
\newcommand{\cours}[1]
{\begin{center} 
  \textit{#1}\\
\end{center}
  }



% \newcommand{\exemple}[1]{\newline~\textbf{Exemple :} #1}
\newcommand{\todo}[1]{\textcolor{red}{#1}}
%\newcommand{\attention}[1]{\newline\faExclamationTriangle~\textbf{Attention :} #1}

% Documentation url (escape \# in the TP document)
\newcommand{\documentation}[1]{\faBookOpen~Documentation : \href{#1}{#1}}

% Clef API
\newcommand{\apikey}[1]{\faKey~Clé API : \lstinline{#1}}
\newcommand{\apiendpoint}[1]{\faGlobe~Url de base de l'API \href{#1}{#1}}

%Listing Python style
\usepackage{color}
\definecolor{slideblue}{RGB}{33,131,189}
\definecolor{green}{RGB}{0,190,100}
\definecolor{blue}{RGB}{121,142,213}
\definecolor{grey}{RGB}{120,120,120}
\definecolor{warning}{RGB}{235,186,1}

\usepackage{listings}
\lstdefinelanguage{texte}{
    keywordstyle=\color{black},
    numbers=none,
    frame=none,
    literate=
           {é}{{\'e}}1
           {è}{{\`e}}1
           {ê}{{\^e}}1
           {à}{{\`a}}1
           {â}{{\^a}}1
           {ù}{{\`u}}1
           {ü}{{\"u}}1
           {î}{{\^i}}1
           {ï}{{\"i}}1
           {ë}{{\"e}}1
           {Ç}{{\,C}}1
           {ç}{{\,c}}1
           {ô}{{\^o}}1,
    columns=fullflexible,keepspaces,
	breaklines=true,
	breakatwhitespace=true,
}
\lstset{
  language=Python,
	%basicstyle=\bfseries\footnotesize,
  basicstyle=\footnotesize\ttfamily\small,
	breaklines=true,
	breakatwhitespace=true,
	commentstyle=\color{grey},
	stringstyle=\color{slideblue},
  keywordstyle=\color{slideblue},
	morekeywords={with, as, True, False, Fraction, Float, join, index, None, main, argparse, self, sort, __eq__, __add__, __mul__ ,__ne__, __radd__, __del__, __ge__, __gt__, split, os, endswith, is_file, scandir, @classmethod, this},
	deletekeywords={id},
	showspaces=false,
	showstringspaces=false,
	columns=fullflexible,keepspaces,
	literate=
           {é}{{\'e}}1
           {è}{{\`e}}1
           {ê}{{\^e}}1
           {à}{{\`a}}1
           {â}{{\^a}}1
           {ù}{{\`u}}1
           {ü}{{\"u}}1
           {î}{{\^i}}1
           {ï}{{\"i}}1
           {ë}{{\"e}}1
           {Ç}{{\,C}}1
           {ç}{{\,c}}1
           {ô}{{\^o}}1,
    numbers=left,
}

\newtcbox{\mybox}{nobeforeafter,colframe=white,colback=slideblue,boxrule=0.5pt,arc=1.5pt, boxsep=0pt,left=2pt,right=2pt,top=2pt,bottom=2pt,tcbox raise base}
\newcommand{\projet}{\mybox{\textcolor{white}{\small projet}}\xspace}
\newcommand{\optionnel}{\mybox{\textcolor{white}{\small Optionnel}}\xspace}
\newcommand{\advanced}{\mybox{\textcolor{white}{\small Pour aller plus loin}}\xspace}
\newcommand{\auto}{\mybox{\textcolor{white}{\small Auto-évaluation}}\xspace}


\usepackage{environ}
\newif\ifShowSolution
\NewEnviron{solution}{
  \ifShowSolution
	\begin{Solutions}{\faTerminal \enskip Solution}
		\BODY
	\end{Solutions}
  \fi}

\usepackage{environ}
\newif\ifShowExample
\NewEnviron{exemple}{
  \ifShowExample
	\begin{Example}{\faTerminal \enskip Exemple}
		\BODY
	\end{Example}
  \fi}

\usepackage{environ}
 \newif\ifShowConseil
  \NewEnviron{conseil}{
    \ifShowConseil
    \begin{Conseils}{\faLightbulb \quad Conseil}
      \BODY
    \end{Conseils}

    \fi}

    \usepackage{environ}
  \newif\ifShowWarning
  \NewEnviron{attention}{
    \ifShowWarning
    \begin{Warning}{\faTriangleExclamation \quad Attention}
      \BODY
    \end{Warning}

    \fi}

  \newif\ifShowNote
  \NewEnviron{note}[1]{
    \ifShowNote
    \begin{Note}{\faEye \quad Comment added by #1 at \date - \currenttime}
      \BODY
    \end{Note}

    \fi}  

    \newif\ifShowDessin
    \NewEnviron{dessin}[1]{
      \ifShowDessin
      \begin{Dessin}{\faPen \quad Figure}
        \BODY
      \end{Dessin}
  
      \fi}  

%\newcommand{\Conseil}[1]{\ifShowIndice\ \newline\faLightbulb[regular]~#1\fi}

\usepackage{array}
\usepackage[T1]{fontenc}
\usepackage{tikz}

\newcolumntype{C}[1]{>{\centering\let\newline\\\arraybackslash\hspace{0pt}}m{#1}}
\setlist[itemize]{label=\textbullet}

\definecolor{DarkBlueContour}{RGB}{0, 0, 150}
\definecolor{LightBlueFill}{RGB}{173, 216, 230}
\tikzset{
    point/.style={circle, draw=DarkBlueContour, fill=LightBlueFill, minimum size=14pt, inner sep=0pt, font=\scriptsize},
    edge/.style={thick, black!70},
    poids/.style={font=\scriptsize, fill=white, inner sep=1pt}
}

\begin{document}
% Change the following values to true to show the solutions or/and the hints
\ShowSolutiontrue
\ShowConseiltrue
\titre
\cours{Algorithmes de graphes - Exercices avancés}

Le but de cette séance est d'appliquer les algorithmes de graphes vus en cours à un problème concret. Dans cette série d'exercices, vous construirez vous-mêmes un graphe pondéré à partir de points du plan, puis vous y chercherez un arbre couvrant minimum (ACM) avec l'algorithme de Kruskal. \\

Le code présenté dans les énoncés se trouve sur Moodle, dans le dossier \lstinline{advanced}.


\section{Relier des points du plan au moindre coût}

Dans toute cette série, on dispose d'une liste \lstinline{points} de $n$ points à coordonnées entières, où \lstinline{points[i]} $= [x_i, y_i]$. Le coût pour relier deux points $[x_i, y_i]$ et $[x_j, y_j]$ est leur \textbf{distance de Manhattan} :
\[
    |x_i - x_j| + |y_i - y_j|
\]
où $|v|$ désigne la valeur absolue de $v$. On cherche le \textbf{coût minimal pour relier tous les points}. Tous les points sont reliés lorsqu'il existe \textbf{exactement un} chemin simple entre n'importe quelle paire de points. \\

\begin{Example}{\faTerminal \quad Exemple}
    \lstinline{points = [[0,0],[1,3],[4,1],[5,5],[2,6]]}. Les points sont numérotés de $P_0$ à $P_4$ selon leur position dans la liste.
    \begin{center}
    \begin{tikzpicture}[scale=0.7]
        \draw[step=1, black!15, very thin] (-0.5,-0.5) grid (6.5,6.5);
        \draw[->, black!50] (-0.5,0) -- (6.7,0) node[right] {\scriptsize $x$};
        \draw[->, black!50] (0,-0.5) -- (0,6.7) node[above] {\scriptsize $y$};
        \node[point] at (0,0) {$P_0$};
        \node[point] at (1,3) {$P_1$};
        \node[point] at (4,1) {$P_2$};
        \node[point] at (5,5) {$P_3$};
        \node[point] at (2,6) {$P_4$};
    \end{tikzpicture}
    \end{center}
    Par exemple, relier $P_1 = [1,3]$ et $P_2 = [4,1]$ coûte $|1-4| + |3-1| = 3 + 2 = 5$.
\end{Example}


% ADVANCED
\begin{Exercice}[15 minutes] [Python] Distance de Manhattan et graphe complet\\

    Le problème ne nous donne pas de graphe : il faut le construire. On considère le \textbf{graphe complet} dont les sommets sont les indices des points ($0, 1, \dots, n-1$) et dans lequel chaque paire de sommets $i$ et $j$ est reliée par une arête non orientée dont le poids est la distance de Manhattan entre \lstinline{points[i]} et \lstinline{points[j]}.

    \begin{enumerate}
        \item Écrivez une fonction \lstinline{distance_manhattan(p, q)} qui retourne la distance de Manhattan entre deux points \lstinline{p} et \lstinline{q}.
        \item Écrivez une fonction \lstinline{construire_aretes(points)} qui retourne la liste de toutes les arêtes du graphe complet, sous la forme de tuples \lstinline{(i, j, w)} avec \lstinline{i < j} et \lstinline{w} le poids de l'arête. Chaque paire de points ne doit apparaître qu'une seule fois.
        \item Combien d'arêtes contient ce graphe pour $n$ points ?
    \end{enumerate}

    \lstinputlisting{code/advanced/aretes_manhattan.py}

    \begin{conseil}
        La fonction \lstinline{abs} de Python retourne la valeur absolue d'un nombre. Pour parcourir chaque paire une seule fois, utilisez deux boucles \lstinline{for} imbriquées et faites commencer la seconde juste après l'indice de la première.
    \end{conseil}

    \begin{solution}
        \lstinputlisting{solutions/advanced/aretes_manhattan.py}

        \textbf{Nombre d'arêtes :} le point $0$ est relié aux $n-1$ points suivants, le point $1$ aux $n-2$ suivants, et ainsi de suite jusqu'au point $n-2$, relié uniquement au point $n-1$. Le graphe contient donc
        \[
            (n-1) + (n-2) + \dots + 1 = \frac{n(n-1)}{2}
        \]
        arêtes, ce qui correspond au nombre d'arêtes d'un graphe complet vu dans la série de base. Pour l'exemple ($n = 5$), on obtient $10$ arêtes. La construction prend donc un temps de $O(n^2)$ et la liste occupe un espace de $O(n^2)$.
    \end{solution}

\end{Exercice}



% ADVANCED
\begin{Exercice}[25 minutes] [Python] Coût minimal pour relier tous les points\\

    Écrivez une fonction \lstinline{cout_minimal(points)} qui retourne le coût minimal pour relier tous les points. On suppose que la liste contient au moins un point. Le fichier fourni contient la classe \lstinline{Graph} de \lstinline{kruskal.py} (série de base), réduite à sa partie Union-Find : \lstinline{find} et \lstinline{union}.\\

    \begin{Example}{\faTerminal \quad Exemples}
        \textbf{Exemple 1}\!: pour \lstinline{points = [[0,0],[1,3],[4,1],[5,5],[2,6]]}, la fonction retourne \lstinline{17}.\\

        \textbf{Exemple 2}\!: pour \lstinline{points = [[-1,2],[3,-2]]}, la fonction retourne \lstinline{8} : il n'y a qu'une seule façon de relier deux points, au coût de $|-1-3| + |2-(-2)| = 8$.\\

        \textbf{Exemple 3}\!: pour \lstinline{points = [[7,-3]]}, la fonction retourne \lstinline{0} : un seul point est déjà relié à lui-même.
    \end{Example}

    \lstinputlisting{code/advanced/cout_minimal.py}

    \begin{conseil}
        Relisez la condition « exactement un chemin simple entre n'importe quelle paire de points ». Quel type de sous-graphe obtient-on ? S'il existe un chemin entre chaque paire de points, le sous-graphe est connexe ; s'il n'en existe qu'un seul, il n'a pas de cycle. Il s'agit donc d'un arbre couvrant, et on cherche celui de poids minimal. Construisez les arêtes comme à la question 1, puis appliquez l'algorithme de Kruskal en additionnant les poids des arêtes retenues.
    \end{conseil}

    \begin{solution}
        Un sous-graphe dans lequel il existe exactement un chemin simple entre chaque paire de sommets est connexe et sans cycle : c'est un \textbf{arbre couvrant}. Le coût minimal pour relier tous les points est donc le poids d'un \textbf{arbre couvrant minimum} du graphe complet de la question 1.

        \lstinputlisting[firstline=27, firstnumber=27, lastline=49]{solutions/advanced/cout_minimal.py}
    \end{solution}
    \begin{solution}
        \textbf{Déroulement sur l'exemple 1.} Les arêtes triées par poids croissant sont :
        $(0,1,4)$, $(1,4,4)$, $(3,4,4)$, $(0,2,5)$, $(1,2,5)$, $(2,3,5)$, $(1,3,6)$, $(2,4,7)$, $(0,4,8)$, $(0,3,10)$.
        \begin{itemize}
            \item $(0,1,4)$, $(1,4,4)$ et $(3,4,4)$ relient chaque fois deux ensembles différents : on les ajoute (coût $4 + 4 + 4 = 12$).
            \item $(0,2,5)$ relie $P_2$, encore isolé, à l'arbre : on l'ajoute (coût $12 + 5 = 17$).
            \item On a $n - 1 = 4$ arêtes : l'arbre couvrant est complet et l'algorithme s'arrête.
        \end{itemize}
        \begin{center}
        \begin{tikzpicture}[scale=0.7]
            \draw[step=1, black!15, very thin] (-0.5,-0.5) grid (6.5,6.5);
            \node[point] (P0) at (0,0) {$P_0$};
            \node[point] (P1) at (1,3) {$P_1$};
            \node[point] (P2) at (4,1) {$P_2$};
            \node[point] (P3) at (5,5) {$P_3$};
            \node[point] (P4) at (2,6) {$P_4$};
            \draw[edge] (P0) -- node[poids] {4} (P1);
            \draw[edge] (P1) -- node[poids] {4} (P4);
            \draw[edge] (P4) -- node[poids] {4} (P3);
            \draw[edge] (P0) -- node[poids] {5} (P2);
        \end{tikzpicture}
        \end{center}

        \textbf{Complexité : $O(n^2 \log n)$.} Construire les $\frac{n(n-1)}{2}$ arêtes prend $O(n^2)$. Les trier prend $O(n^2 \log n^2) = O(n^2 \log n)$, car $\log n^2 = 2 \log n$. Chaque appel à \lstinline{union} est quasiment en temps constant grâce à la compression de chemin et à l'union par rang. Le tri domine donc le temps d'exécution. L'espace utilisé est de $O(n^2)$ pour stocker la liste des arêtes.\\

        \textbf{Remarque :} l'ACM n'est pas forcément unique. Ici, on aurait aussi pu choisir $(1,2,5)$ ou $(2,3,5)$ à la place de $(0,2,5)$ pour relier $P_2$. Le coût minimal, lui, est toujours le même : $17$.
    \end{solution}

\end{Exercice}

\end{document}
